\subsection*{R2-4. Pairwise selection correlation analysis}

\reviewercomment{The pairwise ``selection correlation'' metric is an interesting attempt to quantify lineage-level coherence, but its interpretation remains somewhat unclear.

It appears closely related to co-occurrence patterns, and it is not fully clear whether the observed correlations reflect ecological interactions, shared environmental responses, or methodological effects. We find this analysis potentially insightful, but its conceptual interpretation would benefit from clarification.

We suggest adding a short paragraph linking this metric more explicitly to invasion fitness in a gLV framework. In that context, invasion fitness corresponds to the growth rate of a species when rare in a resident community. The pairwise invasion assays can then be interpreted as empirical approximations of invasion fitness in two-species systems, while coalescence outcomes reflect correlations in invasion success across many species.

This would provide a unifying theoretical framework linking pairwise assays, community coalescence, and the notion of community-level selection.}

We agree that the pairwise selection correlation should be interpreted as origin-correlated invasion success, not as a stand-alone proof of one specific ecological mechanism. We have clarified the metric in Supplementary Note~4: for each coalescence event, the metric asks whether species pairs have concordant post-coalescence fates, either both persisting or both being lost, and compares same-parental-community pairs with cross-parental-community pairs using origin-label permutations. Thus, the metric is related to co-occurrence, but it is applied to a defined perturbation experiment, community coalescence, and tests whether persistence is correlated with parental-community origin rather than whether taxa simply co-occur in an observational sample.

We also added the gLV invasion-fitness connection requested by the reviewer. In the gLV framework, the invasion fitness of species $i$ introduced when rare into a resident community at equilibrium is its per-capita growth rate in that resident background; positive values imply establishment and negative values imply exclusion. This gives a direct mapping between our assays and the coalescence metric: the pairwise invasion assays are empirical two-species approximations of invasion fitness using a 5\% invader frequency, whereas the pairwise selection correlation measures whether many species from the same parental community have correlated invasion outcomes when confronted with another assembled community. To make this quantitative, we added an auxiliary gLV analysis in which excess same-parent invasion concordance, measured relative to a null expectation, increases across the simulated interaction-strength sweep with Pearson $r = 0.870$ and $p = 3.2 \times 10^{-8}$ (Supplementary Fig.~36).

This framing also clarifies the boundary of the inference. A positive pairwise selection correlation shows that species persistence is correlated with parental-community origin, which is the operational signature of community-level selection used in the manuscript. It does not by itself distinguish all possible causal routes, such as direct pairwise competition, shared environmental tolerance, or environmental modification by dominant taxa, and it does not imply that two-species assays fully determine multi-species coalescence outcomes. We therefore now state explicitly in the Discussion that community-level selection is an outcome-level pattern with multiple possible mechanistic routes, and we use the gLV analysis to show that competitive interaction structure is sufficient to generate the pattern, not that it is the only mechanism in the experiments.

\mschange{We revised Results to state that the same-vs-cross selection metric can be interpreted as correlated multi-species invasion success in the gLV framework and to cite the auxiliary invasion-fitness analysis in Supplementary Fig.~36. We added Supplementary Note~4, ``Connection to invasion fitness,'' defining invasion fitness in the rare-invader limit, mapping the pairwise invasion assays and coalescence metric to that framework, and reporting the $\mu$-sweep result ($r = 0.870$, $p = 3.2 \times 10^{-8}$). We also revised the Discussion to define community-level selection operationally as origin-correlated persistence and to state that shared environmental tolerances or environmental filtering can be mechanisms for this outcome rather than alternatives that are excluded by the metric alone.}
